% ============================================================================
%  Chapter 7 — Discussion
%  Thesis: Generative AI, Privacy & Personal Data — TADJ Mohamed Walid, ESI SBA
%
%  DUAL-MODE FILE (standalone OR \input via \ThesisMaster). Sets \label{ch:discussion}.
%  Drafted FROM discussion_notes.md (the banked spine); structure follows it.
%  Measured claims are stated flatly with a \ref to Ch3-6 (where they are proven);
%  synthesis claims that exceed measurement are HEDGED ("we argue", "this suggests",
%  "the most plausible reading") and never asserted as measured. No result is
%  re-derived here. Citations reuse existing references.bib keys (no new entries).
% ============================================================================
\ifdefined\ThesisMaster\else
  \documentclass[11pt,a4paper]{report}
  \usepackage[utf8]{inputenc}
  \usepackage[T1]{fontenc}
  \usepackage{lmodern}
  \usepackage[margin=2.5cm]{geometry}
  \usepackage{booktabs}
  \usepackage{tabularx}
  \usepackage{array}
  \newcolumntype{L}{>{\raggedright\arraybackslash}X}
  \usepackage{amsmath,amssymb}
  \usepackage{microtype}
  \usepackage[hidelinks]{hyperref}
  \begin{document}
\fi

\chapter{Discussion}
\label{ch:discussion}

The preceding chapters established a baseline, built a layered defense, measured its effect
through two instruments, and quantified what it leaves behind. This chapter steps back from the
numbers to ask what they mean together. It does not re-derive any result (each measured claim
points by reference to the chapter that establishes it), and it is careful, throughout, to
keep two kinds of statement apart: what the data shows, and what we argue follows from it. The
discussion's value is in the second kind, but its credibility depends on never dressing the
second as the first.

\section{The methodological headline}
\label{sec:disc-headline}

The central methodological result of this thesis is that measuring a defended
retrieval-augmented system with a single instrument is wrong in a direction the analyst cannot
see. Read through the deterministic rule engine, the full defense stack reduces leakage to
$0.000$; read through the validated judge, it leaves a residual of $17.3\%$
(Chapter~\ref{ch:defenses}). The gap is not measurement noise but the engine's structural
blindness to semantic and cross-document disclosure, recorded as if it were protection. Because
the engine can only credit a leak to a discrete annotated entity on a retrieved document, its
error has a fixed sign: it can only under-count, so its blindness always reads as success
(Chapter~\ref{ch:baseline}). The single-instrument error is therefore systematic and
optimistic, not random.

Here the line between measurement and argument must be drawn. That \emph{this} corpus shows a
$0.000$-versus-$17.3\%$ gap, concentrated in the linkability family, is measured. That the
finding \emph{generalizes} is an argument, and we make it as one: any
annotation-bound or pattern-based detector is structurally unable to see leakage that is
\emph{emergent} (paraphrased, inferred, or re-assembled across documents), because such
leakage corresponds to no discrete entity it can match, and that a defense which suppresses
discrete disclosures while leaving emergent ones intact will therefore read as near-total under
such a detector on any corpus. What travels beyond this study, in other words, is the
\emph{direction} of the error; its \emph{magnitude} (a genuine reduction of roughly forty
percent rather than the near-elimination the engine reports) is specific to this corpus and
should not be read as a prediction.

\section{Five questions, one problem}
\label{sec:disc-throughline}

The five research questions are usually answered in sequence, but they are better
understood as five faces of a single phenomenon, and that reading them this way is what turns a
set of results into a finding. The connecting thread is the nature of re-identification by
linkage. A person reconstructed from individually innocuous attributes scattered across
documents is not a discrete entity sitting in one record; the disclosure is \emph{emergent}, and
this single property explains the otherwise separate results of the four empirical
questions at once.

Consider them in turn, each a face of the same thing. The undefended baseline (RQ1) is
dominated by two families: direct PII extraction at a judged $0.778$ and cross-document
linkability at $0.600$ (Chapter~\ref{ch:baseline}). Retrieval-time redaction (RQ2) is the
decisive layer, cutting the headline from $28.8\%$ to $17.3\%$, but it acts on discrete
entities: direct extraction collapses to $0.222$ while linkability does not move from $0.600$
(Chapter~\ref{ch:defenses}, Chapter~\ref{ch:tradeoff}). The reason is the emergence property:
there is no discrete span for redaction to remove. The same property means the deterministic
engine cannot even \emph{detect} the surviving linkage (its agreement with the human on defended
output is $\kappa = -0.10$, no better than chance; Chapter~\ref{ch:defenses}), so the only
instrument that can see it is the semantic judge, which is precisely the instrument whose
reliability degrades on defended output (RQ3: $\kappa$ falls from $0.827$ to $0.306$), and whose
three missed leaks are \emph{all} re-identification by linkage. And because the insider attacks
by linkage rather than extraction, the insider is the hardest adversary to defend, barely moving
from $0.333$ to $0.278$ while the external and administrator attacks are brought down sharply
(RQ4; Chapter~\ref{ch:tradeoff}). The residual (RQ5) is therefore linkage in both registers, $L$
and $S$ (Chapter~\ref{ch:tradeoff}).

The interpretation we draw from this is the chapter's central claim: \textbf{the leak that is
hardest to prevent is hardest to measure, and not by coincidence.} The same emergence that puts
linkage beyond redaction's reach also puts it beyond a deterministic detector's sight and at the
edge of a semantic judge's reliability. The defense gap and the
measurement gap in this thesis are therefore not two difficulties but one, seen from two sides. One
consequence of this reading deserves emphasis because it cuts against over-claiming: the leaks
the judge missed were exactly the linkage cases the human annotator caught
(Chapter~\ref{ch:defenses}). The residual is thus corroborated by human judgment rather than
inflated by the judge, and the $17.3\%$ figure is best read as a lower bound on what the
instruments can see, not an upper one.

\section{Calibration as a contribution to LLM-as-judge evaluation}
\label{sec:disc-judge}

Because the headline rests on a language-model judge, the thesis treats that judge as an object
of measurement in its own right, and the resulting calibration story is a contribution
independent of the privacy numbers. The measured facts are these: judge--human agreement falls
from $\kappa = 0.827$ on undefended output to $\kappa = 0.306$ on the defended distribution,
with raw agreement (the prevalence-independent companion) falling from $93\%$ to $70\%$; the
error is bidirectional, eleven over-flags (refusals read as disclosures) against three
under-detected linkage leaks; and propagating the two error rates honestly, after the textbook
Rogan--Gladen correction degenerates under the judge-positive-enriched sample, yields a wide
interval of roughly $[5.3\%, 25.8\%]$ that contains $17.3\%$ (Chapter~\ref{ch:defenses}). One
claim is deliberately \emph{not} made: that the true rate lies below $17.3\%$. We treat that as
forbidden, because the calibration sample took every judge-positive, so the visible over-flagging
cannot be projected into a downward correction without selection bias.

Two interpretive lessons follow, and we offer both as arguments. First, on the abandoned
multi-agent panel: an arbitration panel whose members share a model-family lineage
manufactures the appearance of independence without its substance, since an arbitrator from the
same family cannot, even in principle, improve on the judge it arbitrates. This is a concrete
instance of the echo-chamber effect that the surveyed literature anticipated
(\cite{ma2025debate}), and the honest resolution, a single validated judge corroborated across
vendors, is a finding rather than a retreat. Second, this suggests that judge
reliability is \emph{distribution-dependent}: a judge validated on raw output cannot be assumed
reliable on the redacted and refused output a defense produces, and
defended-distribution recalibration should be treated as standard practice rather than an
afterthought. Both lessons answer, in part, the open frontier that \cite{yehudai2025ecosystem}
names, organizational-policy-compliance evaluation, and both are consistent with the poor
metacognitive calibration that \cite{kim2025objexmt} reports for LLM judges generally.

\section{Limitations, and why naming them is the point}
\label{sec:disc-limits}

The limitations of this work are best read not as a checklist but as three clusters that
interact. The first is an \emph{instrument} cluster: a single-provider judge, a same-vendor
language-model redaction probe whose shared blind spot cannot be ruled out, and a
defended-distribution recalibration labeled by a single annotator (the second annotator
covered only the undefended chain). We note that these do
not merely co-exist; they compound, because each is a single-vendor or single-annotator exposure
on the same measurement apparatus. The second is a \emph{corpus} cluster: the synthetic corpus
samples attributes independently, so its absolute numbers do not transfer (only its relative
findings do), its $84$-in-$100$ re-identification figure is an upper bound on intrinsic
uniqueness rather than a conservative estimate, and the corpus-poisoning family rests on a single
scenario (Chapter~\ref{ch:threat}, Chapter~\ref{ch:baseline}). The third is a \emph{utility}
cluster: only three quality dimensions were instrumented, and the cost was measured at the two
extremes rather than isolated at the knee (Chapter~\ref{ch:tradeoff}).

We interpret these characterized limitations as a sign of methodological maturity rather than
weakness, and the point is consistent with the thesis's own headline: if a result is
only as trustworthy as the instrument that measured it, then measuring and bounding the
instrument (the judge's defended-distribution $\kappa$, the noise floor, the upper-bound
reading of the re-identification figure) is not a concession but the discipline the headline
demands. The honest scope of the work follows from the same stance: this is a demonstration of a
measurement methodology on a controlled corpus, not a certification that a deployable product is
safe.

\section{Practical implications}
\label{sec:disc-practical}

For a practitioner deploying enterprise RAG, four implications follow, and we separate the
measured basis of each from the recommendation we draw from it. The deployment recommendation is
the knee (redaction plus conversational refusal): it captures essentially all of the full
stack's measured security ($0.192$ against $0.173$, within the noise floor;
Chapter~\ref{ch:tradeoff}), and the full stack is therefore best understood as
defense-in-depth for attack surfaces the benign workload never exercises, not as a measured
improvement. Redaction is the decisive layer, the only one that clears the noise
floor (Chapter~\ref{ch:defenses}), and budget should reflect that. On the residual, linkage
will not yield to redaction and requires a different class of defense entirely:
roughly two-thirds of the residual is a detector-recall gap that better recognizers could close,
but the remaining third is genuine re-identification that needs entity-resolution or
directory-aware controls that redaction cannot provide (Chapter~\ref{ch:tradeoff}). Finally, and
most directly tied to the headline, a defended system should be evaluated with a
validated judge recalibrated on defended output, not a detector alone: a deployment that trusts
the rule engine will believe it has eliminated leakage when it has reduced it by about forty
percent.

\section{From measurement to practice: recommendations and regulatory bearing}
\label{sec:disc-recommendations}

The practical implications above translate into guidance for practitioners under the one discipline that
holds the rest of this thesis: every recommendation below traces to a specific measured result, and none
is stated more strongly than the evidence supports. The findings show a layered defense that reduces
leakage by about forty percent but does not eliminate it, a linkage residual that resisted three separate
mechanisms, and instruments that disagree on defended output; the recommendations reflect exactly
that, no more.

\begin{enumerate}[leftmargin=1.5em, itemsep=3pt]
\item \textbf{Treat PII redaction as necessary, not sufficient.} Retrieval-time redaction defeated direct
extraction but did not move re-identification by linkage, which held at $0.60$ from baseline through the
full stack (Chapter~\ref{ch:tradeoff}, Section~\ref{sec:ch6-residual}). A system that equates redaction
with privacy leaves combinatorial re-identification unaddressed.
\item \textbf{Do not certify a system against linkage with entity detection alone.} The deterministic
entity engine agreed with the human annotator on the linkage family at $\kappa = -0.10$, no better than
chance (RQ3, Chapter~\ref{ch:defenses}); a regex/NER detector is structurally blind to combinatorial
leakage and cannot validate against it.
\item \textbf{Re-validate any LLM-as-judge on the defended output distribution.} The validated judge's
agreement with the human fell from $\kappa = 0.827$ on undefended output to $\kappa = 0.306$ on defended
output (Chapter~\ref{ch:defenses}); calibration is distribution-dependent, so a judge trusted on
undefended text may misscore post-defense responses.
\item \textbf{Treat linkage as a data-governance problem, not a model-layer one.} Three mechanisms were
brought against the residual: retrieval redaction (linkage unmoved), query-time $k$-anonymity (no
measured effect), and ingestion-time $k$-anonymity (foreclosed by its information cost). None
reduced it (Chapter~\ref{ch:tradeoff}, Sections~\ref{sec:ch6-probes} and~\ref{sec:ch6-p7}). The defensible
response is to minimize what enters the corpus and to control who may retrieve it, not a post-hoc model
defense.
\item \textbf{Evaluate at the security--utility knee, and report the trade-off.} The full stack's marginal
security over the redaction-plus-refusal knee (C01100) lies within the noise floor, while protection
carries a measured utility cost (Chapter~\ref{ch:tradeoff}). Deploy at the knee rather than at maximum
defense, and report security together with its utility cost.
\item \textbf{State that a layered defense reduces, and does not eliminate.} The measured reduction,
the headline $28.8\%$ to $17.3\%$, about forty percent, is carried by redaction with refusal at the
knee (Chapter~\ref{ch:defenses}, Chapter~\ref{ch:tradeoff}) and does not reach zero. Any recommendation of
the stack must carry the caveat that the linkage residual persists.
\end{enumerate}

These findings bear on obligations under the instruments this work aligns with: the GDPR
(Regulation (EU) 2016/679), the NIS2 Directive (Directive (EU) 2022/2555), and ISO/IEC 27001:2022.
What follows, though, is a \emph{mapping of measured results to regulatory themes, not a legal compliance
assessment}. The linkage residual and the re-identification finding (four quasi-identifiers already make
$84$ of $100$ employees unique, Chapter~\ref{ch:formal}) bear on the GDPR's data-minimization principle
(Art.\ 5(1)(c)) and its standard for anonymization: a corpus carrying these quasi-identifiers supports
re-identification and is, on this evidence, not anonymous in that sense, which grounds the
corpus-minimization recommendation. The access-control recommendation bears on the security-of-processing
and access-control controls common to the GDPR (Art.\ 32), NIS2, and ISO/IEC 27001. The judge-calibration
finding bears on the accountability principle (Art.\ 5(2)): automated privacy-evaluation evidence is only
as valid as the instrument's calibration on the distribution it scores. We claim no more than relevance to
these obligations; we do not certify compliance, and the thesis is not a legal document. \Cref{tab:disc-recommendations} traces each recommendation to the measured finding behind it and the regulatory theme it bears on.

\begin{table}[h]
\centering
\caption[Traceability of recommendations to findings and regulatory themes]{Traceability of each recommendation to its measured finding and the regulatory theme it bears on.
Each recommendation cites a measured result; the regulatory column is a mapping, not a compliance claim.}
\label{tab:disc-recommendations}
\begin{tabularx}{\textwidth}{LLL}
\toprule
\textbf{Recommendation} & \textbf{Grounding finding (chapter)} & \textbf{Regulatory bearing}\\
\midrule
Redaction is necessary, not sufficient & Linkage unmoved at $0.60$ (\S\ref{sec:ch6-residual}) & GDPR data minimization (Art.\ 5(1)(c))\\
\midrule
Entity detection cannot certify linkage & Engine $\kappa=-0.10$, RQ3 (Ch.~\ref{ch:defenses}) & GDPR accountability (Art.\ 5(2))\\
\midrule
Recalibrate the judge on defended output & Judge $\kappa$ $0.827\!\to\!0.306$ (Ch.~\ref{ch:defenses}) & GDPR accountability (Art.\ 5(2))\\
\midrule
Linkage is data-governance, not model-layer & Three mechanisms failed (\S\ref{sec:ch6-probes},~\ref{sec:ch6-p7}) & GDPR data min.; NIS2 / ISO 27001 access control\\
\midrule
Deploy at the security--utility knee & Knee $\approx$ full stack within noise (Ch.~\ref{ch:tradeoff}) & GDPR proportionate measures (Art.\ 25)\\
\midrule
Layered defense reduces ($\sim$40\%), not eliminates & $28.8\%\!\to\!17.3\%$ at C01100 (Ch.~\ref{ch:defenses},~\ref{ch:tradeoff}) & GDPR security of processing (Art.\ 32)\\
\bottomrule
\end{tabularx}
\end{table}

The honest position is the strength of this section: these are implications of the thesis's measurements
for practice (``our measurements imply X''), not a recipe for a guaranteed-private system, which the
thesis did not build or certify. Each recommendation stays within what was measured, and the regulatory
mapping identifies which obligations the findings touch without asserting compliance.

\subsection{Governance controls for the linkage residual}
\label{sec:disc-governance}

The three negatives of the preceding chapters (redaction, the query-time gate L7, and the
ingestion-time generalization P7) establish that linkage cannot be defended at the model layer
without destroying utility, because re-identification is a property of the information itself
\cite{sweeney2002kanonymity}: four quasi-identifiers make $84$ of $100$ employees unique
(Chapter~\ref{ch:formal}). What follows is therefore \emph{not} another model-layer defense, and
\emph{not} a tested defense stack with a measured F2; none of these controls was run as an experiment.
It is a reasoned proposal, grounded in the thesis's own findings and in the established literature, for the
governance-of-access direction those negatives point to. The controls (\Cref{fig:ch7-governance}) change the problem (what is
ingested, who may retrieve, and how much may be enumerated), and each is stated with the boundary of what
it does and does not cover. Only one case, the external attacker, earns the word \emph{solves}; everywhere
else the controls \emph{bound}, \emph{reduce}, or \emph{make accountable}.

\begin{figure}[h]
\centering
\resizebox{\textwidth}{!}{%
\begin{tikzpicture}[font=\small, >={Latex[length=2.4mm]},
  stage/.style={draw=figslate!80!black, fill=figslate!12, rounded corners=3pt, minimum height=11mm, minimum width=20mm, align=center, thick},
  c1/.style={draw=figteal!80!black, fill=figteal!16, rounded corners=3pt, minimum height=14mm, minimum width=26mm, align=center, very thick, text=figteal!30!black},
  c2/.style={draw=figblue!80!black, fill=figblue!14, rounded corners=3pt, minimum height=14mm, minimum width=26mm, align=center, very thick, text=figblue!40!black},
  c3/.style={draw=figamber!85!black, fill=figamber!20, rounded corners=3pt, minimum height=14mm, minimum width=26mm, align=center, very thick, text=figorange!55!black},
  ar/.style={->, thick, figslate!70!black}]
\node[stage] (corpus) {Corpus};
\node[c1, right=5mm of corpus] (c1) {\textbf{C1}\\\footnotesize data\\minimization};
\node[stage, right=5mm of c1] (ret) {Retrieval};
\node[c2, right=5mm of ret] (c2) {\textbf{C2}\\\footnotesize access\\control};
\node[stage, right=5mm of c2] (gen) {Generation};
\node[c3, right=5mm of gen] (c3) {\textbf{C3}\\\footnotesize aggregation\\control};
\node[stage, right=5mm of c3] (ans) {Answer};
\draw[ar](corpus)--(c1);\draw[ar](c1)--(ret);\draw[ar](ret)--(c2);\draw[ar](c2)--(gen);\draw[ar](gen)--(c3);\draw[ar](c3)--(ans);
\node[draw=figgreen!75!black, fill=figgreen!15, rounded corners=3pt, align=center, text width=48mm, font=\footnotesize, text=figgreen!32!black, below=11mm of ret] (ext) {\textbf{External attacker:} \emph{solved}\\(no sensitive tier served)};
\node[draw=figcoral!75!black, fill=figcoral!12, rounded corners=3pt, align=center, text width=54mm, font=\footnotesize, text=figcoral!45!black, below=11mm of gen] (auth) {\textbf{Authorized user $+$ auxiliary knowledge:}\\\emph{bounded}, not eliminated};
\end{tikzpicture}}
\caption[The governance-of-access controls]{The three governance-of-access controls as a
defense-in-depth: \textbf{C1} minimize what enters the corpus, \textbf{C2} gate retrieval by
authorization, \textbf{C3} control aggregation at the answer layer. Against the external,
unauthorized attacker the approach \emph{solves} the problem (nothing sensitive is served); for
the authorized user who brings outside knowledge the residual is \emph{bounded} --- minimized and
made accountable --- but not eliminated.}
\label{fig:ch7-governance}
\end{figure}

\paragraph{Control 1: purpose-bound data minimization at ingestion.} Index only the fields the system's
defined purpose requires, and never ingest the quasi-identifiers it does not. This is confident because it
is definitional: absent data cannot leak, and, unlike P7, which generalized \emph{all} fields and so
gutted the corpus (Section~\ref{sec:ch6-p7}), it removes only \emph{un-needed} fields, costing no needed
utility. Its boundary is equally clear: it eliminates linkage only if the fields the purpose \emph{does}
require are themselves non-identifying, which for an HR corpus they are not. It therefore \emph{reduces} the
linkage surface to the purpose-minimum; it does not clear the residual.

\paragraph{Control 2: authorization-gated retrieval.} Partition the corpus by sensitivity and enforce
role- or attribute-based access control at the retrieval layer, so a request is served only the content its
identity is authorized for \cite{ozgur2024rbacrag}. This case splits, and the split is the whole point.
\emph{For the external, unauthorized attacker} (profile P2, whose authorization threshold is zero, so any
disclosure is a leak), serving no sensitive content \emph{solves} the problem: with nothing sensitive
retrieved, there is nothing to link. This is the one place the word is earned: it is definitional and
tied directly to the threat model. \emph{For an authorized, internal user} (P1) it does \emph{not} stop
linkage, and we do not claim it does: linkage rides on the low- and medium-severity quasi-identifiers
(department, age, office, gender) that clearance-gating passes through, and the re-identification analysis
shows roughly four such attributes (several of severity $\le 2$) make $84$ of $100$ employees unique
(Chapter~\ref{ch:formal}), with the age--gender combination precisely the gap the detector misses
(Chapter~\ref{ch:tradeoff}). An authorized user, particularly one who already
knows a few attributes of a target from outside the system, clears the access check and still completes
the re-identification from the quasi-identifiers the system legitimately returns. This within-clearance,
auxiliary-knowledge linkage is bounded, narrowed by minimization (Control~1) and made accountable by
audit, but it is not eliminated. We note
honestly, too, that the entity-level access control this control assumes is \textsc{untested}: the
document-level access control the thesis did evaluate proved too coarse, because every legitimate document
carries some high-severity content, so filtering by clearance blocked all retrieval
(Chapter~\ref{ch:defenses}); the finer-grained version is a plausible refinement, not a measured result.

\paragraph{Control 3: aggregation control at the answer layer.} The linkage attack is a request to
narrow a candidate set; the classical statistical-disclosure-control response is to answer such queries with
suppressed aggregates (counts rather than individuals, with sets smaller than a threshold withheld),
backed by a per-session query budget. The suppression \emph{mechanism} is well established
\cite{adam1989sdc}. Two boundaries are honest, however. First, it depends on \emph{detecting} that a query is a set-narrowing one, and
detection is exactly what this thesis found hard: the deterministic engine's agreement with the human on
linkage is $\kappa = -0.10$ (Chapter~\ref{ch:defenses}), no better than chance, so the confidence lives
in the suppression, not in the detection. Second, suppression is not unconditional: the disclosure-control
and differential-privacy literature shows that small suppressed cells remain recoverable and that unlimited
accurate queries eventually leak, which is why the per-session budget is required
\cite{domingoferrer2021dplimits}. It is a strong mitigation for the pattern, not a guarantee.

\paragraph{Synthesis.} The honest direction is therefore a defense-in-depth of \emph{governance} controls,
not a model-layer fix: minimize what enters the corpus (Control~1), gate retrieval by authorization
(Control~2), and throttle or aggregate enumeration while auditing the rest (Control~3). For the unauthorized
attacker this is a definitional solve; for the authorized user the within-clearance residual is
\emph{bounded}, minimized, throttled, and made accountable, but not eliminated, and the thesis is
right that it cannot be eliminated cheaply. These controls are a reasoned, finding- and literature-grounded
proposal, not a tested stack; turning them into measured defenses is future work.

\paragraph{A foreclosed extension: cross-query disclosure accounting.} One further mechanism deserves a
line, as the session-level companion to Control~3, because it can be foreclosed without spending: a
per-session \emph{disclosure accountant} that tracks, across a user's query sequence, the intersection of
the candidate sets the session has narrowed to, and throttles once that running set falls below $k$. It is
foreclosed as a standalone defense by a finding this thesis already recorded: the linkage here is
\emph{single-chunk}, with the identifying combination present the moment a record is retrieved, not
accumulated across turns (Section~\ref{sec:ch6-probes}). An accountant therefore fires on the \emph{first}
query, where the candidate set already collapses below $k$: single-query disclosure, out of scope for a
mechanism whose only purpose is to bound \emph{cross-query} narrowing. Its one honest home is
\emph{contingent}: cross-query accounting bounds linkage only when layered on top of single-query
(combined-constraint) blocking, catching an attacker who, denied the one-shot query, tries to reconstruct
it incrementally. It is therefore a \emph{complement to, not a substitute for}, single-query control, not a
standalone defense.

\section{Future work}
\label{sec:disc-future}

Each direction for future work answers a limitation named above. The single-vendor/single-annotator
exposure of the instrument cluster calls for a second human annotator and a full multi-vendor
judge panel, and for treating defended-distribution recalibration (lifting the judge above the
measured $\kappa = 0.306$) as a first-class task. The corpus cluster calls for validation
beyond the single Enron probe, where the open question is the judge-transfer over-flagging
observed on real prose. The residual itself calls for the linkage-specific defenses the practical
implications identify: entity-resolution controls for the genuine re-identification third, and
per-attribute recognizers for the detector-recall two-thirds. A query-time $k$-anonymity gate over
the quasi-identifiers was tested directly (Section~\ref{sec:ch6-probes}) and did not move the
residual, which points this work toward \emph{ingestion-time} generalization (breaking the
identifying combination before it enters a retrievable unit) rather than a query-time filter. The utility cluster calls for
per-configuration measurement that isolates the knee's marginal cost. And because every attacker
profile here is analyzed in isolation, compositional, multi-actor threats, where a
supply-chain compromise hands an external adversary privileged access, remain open, and we
expect them to be the setting in which the linkage residual matters most.

\ifdefined\ThesisMaster\else
  \bibliographystyle{plain}
  \bibliography{references}
  \end{document}
\fi
